\section{Proof of Theorem 1: Exponential Mean-Bias Amplification}

\subsection{Setup}

Let $x_1,\dots,x_N\in\mathbb{R}^D$ be token representations and write
\[
 x_j = \mu_{c(j)} + \varepsilon_j,
\qquad j=1,\dots,N,
\]
with clusters $C_m$, centroids $\mu_m$, and mixture mean
\[
 \mu = \sum_m w_m \mu_m,
\qquad w_m = \frac{|C_m|}{N},
\qquad \mu\neq 0.
\]
Define the global means
\[
 \bar x := \frac{1}{N}\sum_{j=1}^N x_j,
\qquad
 \bar y := \frac{1}{N}\sum_{i=1}^N y_i.
\]

Single-head attention is
\[
 a_{ij} = \frac{\exp(\langle x_i,x_j\rangle/\tau)}{\sum_k \exp(\langle x_i,x_k\rangle/\tau)},
\qquad
 y_i = \sum_j a_{ij} x_j.
\]

Define mean-energy ratios
\[
 R_X^2 = \frac{\|\bar x\|^2}{\frac1N\sum_i \|x_i\|^2},
\qquad
 R_Y^2 = \frac{\|\bar y\|^2}{\frac1N\sum_i \|y_i\|^2}.
\]

\subsection{Assumptions}

\paragraph{A1 (Zero-mean intra-cluster residuals).}
For each cluster $C_m$,
\[
 \frac{1}{|C_m|}\sum_{j\in C_m} \varepsilon_j = 0.
\]

\paragraph{A2 (Softmax logit gap).}
There exists $\Delta>0$ such that for all $i\in C_m$, $j\in C_m$, and $j'\notin C_m$,
\[
 \langle x_i,x_j\rangle - \langle x_i,x_{j'}\rangle \ge \Delta.
\]

\paragraph{A3 (Residual energy).}
Define
\[
 \sigma^2 := \frac{1}{N}\sum_{i=1}^N \|\varepsilon_i\|^2.
\]

\paragraph{A4 (Within-cluster spread).}
For each $i\in C_m$, the in-cluster weights are not degenerate:
\[
 \sum_{j\in C_m} a_{ij}^2 \le \frac{A_{\mathrm{in}}(i)^2}{|C_m|} + (1-A_{\mathrm{in}}(i))^2.
\]
(Equality holds for uniform weights within the cluster.)

\subsection{Mean and Input Energy}

By A1, the empirical mean equals the mixture mean:
\[
 \bar x = \mu.
\]
The input energy decomposes as
\[
 \frac1N\sum_i \|x_i\|^2 = \|\mu\|^2 + \sigma^2,
\]
so
\[
 R_X^2 = \frac{\|\mu\|^2}{\|\mu\|^2 + \sigma^2}.
\]

\subsection{Exponential Suppression of Cross-Cluster Mass}

Fix $i\in C_m$ and define in-cluster mass
\[
 A_{\mathrm{in}}(i) := \sum_{j\in C_m} a_{ij}.
\]
From A2,
\[
 \frac{a_{ij'}}{a_{ij}} \le \exp\!\left(-\frac{\Delta}{\tau}\right),
\qquad j\in C_m,\ j'\notin C_m,
\]
which yields
\[
 1-A_{\mathrm{in}}(i)
 \le \frac{N-|C_m|}{|C_m|}\,e^{-\Delta/\tau}.
\]
Define
\[
 \eta := \max_i (1-A_{\mathrm{in}}(i))
 \le \max_m \frac{N-|C_m|}{|C_m|}\,e^{-\Delta/\tau}.
\]

\subsection{Variance Contraction}

Write
\[
 y_i = \sum_j a_{ij}\mu_{c(j)} + \sum_j a_{ij}\varepsilon_j.
\]
Let $\tilde\varepsilon_i := \sum_j a_{ij}\varepsilon_j$. By convexity,
\[
 \|\tilde\varepsilon_i\|^2 \le \sum_j a_{ij}\|\varepsilon_j\|^2.
\]
Using A4 and splitting weights into in-cluster and out-of-cluster parts,
\[
 \|\tilde\varepsilon_i\|^2
 \le
 \Big(\sum_{j\in C_m} a_{ij}^2 + (1-A_{\mathrm{in}}(i))\Big)\,\max_j \|\varepsilon_j\|^2.
\]
Averaging over $i$ and using $\sum_j a_{ij}^2 \le A_{\mathrm{in}}(i)^2/|C_m| + (1-A_{\mathrm{in}}(i))^2$,
we obtain a contraction of residual energy by a factor
\[
 1-\kappa,\qquad
 \kappa := c_0\,(1-\eta),
\]
for some constant $c_0>0$ depending only on the minimum cluster size.
Thus
\[
 \frac1N\sum_i \|\tilde\varepsilon_i\|^2 \le (1-\kappa)\sigma^2.
\]
Since $\eta=O(e^{-\Delta/\tau})$, we have
\[
 \kappa = \Theta(1-e^{-\Delta/\tau}).
\]

\subsection{Output Energy and Amplification}

The output mean satisfies
\[
 \bar y = \mu + O(e^{-\Delta/\tau}),
\]
and the output energy obeys
\[
 \frac1N\sum_i \|y_i\|^2 = \|\mu\|^2 + (1-\kappa)\sigma^2 + O(e^{-\Delta/\tau}).
\]
Therefore
\[
 R_Y^2
 =
 \frac{\|\mu\|^2}{\|\mu\|^2+(1-\kappa)\sigma^2}
 + O(e^{-\Delta/\tau}).
\]
Subtracting $R_X^2$ gives
\[
 R_Y^2 - R_X^2
 =
 \frac{\|\mu\|^2\,\kappa\,\sigma^2}{(\|\mu\|^2+\sigma^2)(\|\mu\|^2+(1-\kappa)\sigma^2)}
 + O(e^{-\Delta/\tau}).
\]
Since $\kappa=\Theta(1-e^{-\Delta/\tau})>0$ for $\Delta>0$, we obtain the stated
lower bound with a constant $c>0$ depending only on cluster balance.\qed